\section{Arquitetura Geral}
A arquitetura implementada para o sistema de suporte analítico assenta num processo de integração que reúne dados provenientes de diferentes fontes, aplica transformações adequadas e prepara a sua posterior carga no repositório analítico. Esta secção apresenta a visão geral do fluxo de dados e dos componentes envolvidos, permitindo compreender como a informação é organizada antes de ser disponibilizada no armazém de dados.

A Figura \ref{fig:arquiteturaetl} representa a arquitetura do processo ELT utilizado no sistema analítico. No lado esquerdo encontram se as fontes de dados. Estas incluem ficheiros fornecidos pelo ICNF, dados meteorológicos em formato bruto e informação proveniente da internet como a cobertura mediática. Os dados extraídos seguem para a área de preparação onde ocorre a transformação. Esta área inclui estruturas de suporte, um espaço de quarentena para validação e a camada de transformação que organiza os dados em zonas específicas de integração. Após a preparação, a informação é enviada para a área de carregamento onde é finalmente armazenada no data warehouse.

\begin{figure}
	\centering
	\includegraphics[width=1.0\linewidth]{imagens/Arquitetura_ETL}
	\caption{Arquitetura ELT (Estrutura da DSA)}
	\label{fig:arquiteturaetl}
\end{figure}
